\clearpage
\section{Design}

The implementation uses a token-queue protocol on RabbitMQ.
\begin{itemize}
    \item Each critical section uses its own durable token queue.
    \item The queue contains exactly one persistent token when the section is available.
    \item \texttt{enter()} consumes the token with manual acknowledgments and blocks until it is delivered.
    \item \texttt{exit()} cancels the local consumer and requeues the same delivery with \texttt{basicNack(..., requeue = true)}.
    \item A temporary exclusive bootstrap-lock queue serializes initialization and keeps the protocol recoverable.
\end{itemize}

This design exposes a small public API: \texttt{enter()},
\texttt{exit()}, and \texttt{close()}. Recovery is broker-driven, so if a
process crashes while holding the token, RabbitMQ requeues the unacknowledged
delivery automatically.

The bootstrap lock is important because it removes the only race that could
otherwise create duplicate tokens. Once the token exists, the queue itself
becomes the single source of truth:
the presence of the token means the section is free, and the absence of the
token means somebody owns it.

\subsection{Trade-offs}

The middleware deliberately solves one problem well instead of becoming a
generic coordination framework. It models one critical section as one token
queue, which keeps the API tiny and the reasoning straightforward. The
trade-off is that the broker is part of the availability story, but that is
acceptable here because the assignment explicitly asks for a MOM-based
solution.
